% Options for packages loaded elsewhere
\PassOptionsToPackage{unicode}{hyperref}
\PassOptionsToPackage{hyphens}{url}
\documentclass[
]{article}
\usepackage{xcolor}
\usepackage[margin=1in]{geometry}
\usepackage{amsmath,amssymb}
\setcounter{secnumdepth}{-\maxdimen} % remove section numbering
\usepackage{iftex}
\ifPDFTeX
  \usepackage[T1]{fontenc}
  \usepackage[utf8]{inputenc}
  \usepackage{textcomp} % provide euro and other symbols
\else % if luatex or xetex
  \usepackage{unicode-math} % this also loads fontspec
  \defaultfontfeatures{Scale=MatchLowercase}
  \defaultfontfeatures[\rmfamily]{Ligatures=TeX,Scale=1}
\fi
\usepackage{lmodern}
\ifPDFTeX\else
  % xetex/luatex font selection
\fi
% Use upquote if available, for straight quotes in verbatim environments
\IfFileExists{upquote.sty}{\usepackage{upquote}}{}
\IfFileExists{microtype.sty}{% use microtype if available
  \usepackage[]{microtype}
  \UseMicrotypeSet[protrusion]{basicmath} % disable protrusion for tt fonts
}{}
\makeatletter
\@ifundefined{KOMAClassName}{% if non-KOMA class
  \IfFileExists{parskip.sty}{%
    \usepackage{parskip}
  }{% else
    \setlength{\parindent}{0pt}
    \setlength{\parskip}{6pt plus 2pt minus 1pt}}
}{% if KOMA class
  \KOMAoptions{parskip=half}}
\makeatother
\usepackage{graphicx}
\makeatletter
\newsavebox\pandoc@box
\newcommand*\pandocbounded[1]{% scales image to fit in text height/width
  \sbox\pandoc@box{#1}%
  \Gscale@div\@tempa{\textheight}{\dimexpr\ht\pandoc@box+\dp\pandoc@box\relax}%
  \Gscale@div\@tempb{\linewidth}{\wd\pandoc@box}%
  \ifdim\@tempb\p@<\@tempa\p@\let\@tempa\@tempb\fi% select the smaller of both
  \ifdim\@tempa\p@<\p@\scalebox{\@tempa}{\usebox\pandoc@box}%
  \else\usebox{\pandoc@box}%
  \fi%
}
% Set default figure placement to htbp
\def\fps@figure{htbp}
\makeatother
\setlength{\emergencystretch}{3em} % prevent overfull lines
\providecommand{\tightlist}{%
  \setlength{\itemsep}{0pt}\setlength{\parskip}{0pt}}
\usepackage{bookmark}
\IfFileExists{xurl.sty}{\usepackage{xurl}}{} % add URL line breaks if available
\urlstyle{same}
\hypersetup{
  pdftitle={RMarkdown Lesson 2},
  hidelinks,
  pdfcreator={LaTeX via pandoc}}

\title{RMarkdown Lesson 2}
\author{}
\date{\vspace{-2.5em}}

\begin{document}
\maketitle

\section{How It Works}\label{how-it-works}

This is an R Markdown file, a plain text file that has the extension
\texttt{.Rmd}. You can open a copy
\href{https://posit.cloud/content/181856}{here} on RStudio Cloud.

\pandocbounded{\includegraphics[keepaspectratio]{https://d33wubrfki0l68.cloudfront.net/aa8ef57652bb9b7d6aa18f6c38b6b45a510a8960/1d9d3/lesson-images/how-1-file.png}}

Notice that the file contains three types of content:

\begin{itemize}
\item
  An (optional) YAML header surrounded by \texttt{-\/-\/-}s
\item
  R code chunks surrounded by
  \texttt{\textasciigrave{}\textasciigrave{}\textasciigrave{}}s
\item
  text mixed with simple text formatting
\end{itemize}

\subsection{A Notebook Interface}\label{a-notebook-interface}

When you open the file in the RStudio IDE, it becomes a
\href{https://bookdown.org/yihui/rmarkdown/notebook.html}{notebook
interface for R}. You can run each code chunk by clicking the
\pandocbounded{\includegraphics[keepaspectratio]{https://d33wubrfki0l68.cloudfront.net/18153fb9953057ee5cff086122bd26f9cee8fe93/3aba9/images/notebook-run-chunk.png}}
icon. RStudio executes the code and display the results inline with your
file.

\pandocbounded{\includegraphics[keepaspectratio]{https://d33wubrfki0l68.cloudfront.net/07a00dd9669405f3cba06ef333db180295466252/7b153/lesson-images/how-2-chunk.png}}

\subsection{Rendering output}\label{rendering-output}

To generate a report from the file, run the \texttt{render} command:

\begin{verbatim}
library(rmarkdown) 
render("1-example.Rmd")
\end{verbatim}

Better still, use the ``Knit'' button in the RStudio IDE to render the
file and preview the output with a single click or keyboard shortcut
(⇧⌘K).

\pandocbounded{\includegraphics[keepaspectratio]{https://d33wubrfki0l68.cloudfront.net/96ec0c54c6d64ea2ec3665db9b3b781962ff6339/5cee1/lesson-images/how-3-output.png}}

R Markdown generates a new file that contains selected text, code, and
results from the .Rmd file. The new file can be a finished
\href{https://bookdown.org/yihui/rmarkdown/html-document.html}{web
page},
\href{https://bookdown.org/yihui/rmarkdown/pdf-document.html}{PDF},
\href{https://bookdown.org/yihui/rmarkdown/word-document.html}{MS Word}
document,
\href{https://bookdown.org/yihui/rmarkdown/ioslides-presentation.html}{slide
show},
\href{https://bookdown.org/yihui/rmarkdown/notebook.html}{notebook},
\href{https://bookdown.org/yihui/rmarkdown/tufte-handouts.html}{handout},
\href{https://bookdown.org/}{book},
\href{https://rmarkdown.rstudio.com/flexdashboard/index.html}{dashboard},
\href{https://bookdown.org/yihui/rmarkdown/r-package-vignette.html}{package
vignette} or \href{https://rmarkdown.rstudio.com/formats.html}{other
format}.

\subsection{How it works}\label{how-it-works-1}

When you run \texttt{render}, R Markdown feeds the .Rmd file to
\href{http://yihui.name/knitr/}{knitr}, which executes all of the code
chunks and creates a new markdown (.md) document which includes the code
and its output.

\pandocbounded{\includegraphics[keepaspectratio]{https://d33wubrfki0l68.cloudfront.net/61d189fd9cdf955058415d3e1b28dd60e1bd7c9b/b739c/lesson-images/rmarkdownflow.png}}

The markdown file generated by knitr is then processed by
\href{http://pandoc.org/}{pandoc} which is responsible for creating the
finished format.

This may sound complicated, but R Markdown makes it extremely simple by
encapsulating all of the above processing into a single \texttt{render}
function.

\end{document}
